% Generated by report/tools/tables.py from the result files of the repository. Do not edit.
\begin{table}[tbp]
\centering
\footnotesize
\setlength{\tabcolsep}{3pt}
\caption[Robustesse des verdicts du sondage]{Contrôles de robustesse du sondage, tous exploratoires (écarts 21, 27 et 29 de \file{study/PROTOCOL.md}) : ils disent si un verdict est fragile, sans en changer aucun. Colonnes : les quatre critères qui ont changé (Q5, Y9, X7, X2) et les deux dont l'intervalle évite zéro sans passer Holm (Q6, Y12). Différences en points (RL classique 2026 moins notre lecture de 2018), intervalle de Newcombe à 95~\%, test de Fisher, Holm sur la famille de la grille du critère comme dans la comparaison principale. (a1) le RL classique de 2026 coupé selon sa provenance (premier tirage sans tri sur le titre, file A3 avec tri), puis $p$ de Fisher entre les deux provenances au sein de 2026 ; (a2) $p$ de Fisher entre articles de 2026 lus sur arXiv (45) et sur le site d'une conférence (5), une coupe trop déséquilibrée pour trancher (écart 21 c) ; (a3) test de Cochran-Mantel-Haenszel stratifié par conférence (correction de Haldane) ; (a4) comparaison refaite sans chaque lecteur principal tour à tour : étendue des différences et plus basse borne inférieure ; un critère est fragile quand la différence change de signe ou que l'intervalle contient zéro avec au moins 35 articles de chaque côté. Partie basse : sans les treize articles lus en partie par \texttt{claude-opus-5} (écart 29) ; avec les 10 décisions « non » qu'un arbitre à l'aveugle a retournées après la recherche par mots-clés, puis avec aussi les 24 décisions jugées douteuses, sur 377 décisions « non » examinées (écart 27). \textbf{En gras}, les écarts dont l'intervalle exclut zéro et les $p$ sous 0,05 (sauf les $p$ d'homogénéité de (a1) et (a2), où un $p$ élevé est attendu). La thèse « écoutée là où cela coûtait peu » est rejetée par le test à l'aveugle (b) : corrélation de Spearman entre la différence observée et le coût médian noté par trois évaluateurs +0,16 ($p$ par permutation 0,52), avec le risque −0,04 ($p$ = 0,89). Sources : \file{study/robustness.json}, \file{study/extension\_stats\_sans\_opus5.json}, \file{study/extension\_stats\_falsify\_turned.json}, \file{study/extension\_stats\_falsify\_bound.json}, \file{study/falsify\_no.json}.}
\label{tab:robustesse}
\begin{tabularx}{\linewidth}{@{}>{\raggedright\arraybackslash}Xcccccc@{}}
\toprule
Contrôle & Q5 & Y9 & X7 & X2 & Q6 & Y12 \\
\midrule
Comparaison principale : différence [IC] & \textbf{\makecell{+43\\{}[+23\,;\,+58]}} & \textbf{\makecell{+36\\{}[+17\,;\,+52]}} & \textbf{\makecell{+31\\{}[+11\,;\,+47]}} & \textbf{\makecell{+26\\{}[+8\,;\,+41]}} & \textbf{\makecell{+22\\{}[+3\,;\,+39]}} & \textbf{\makecell{+12\\{}[+1\,;\,+24]}} \\
\quad $p$ Holm & \textbf{0,0004} & \textbf{0,004} & \textbf{0,02} & \textbf{0,04} & 0,24 & 0,66 \\
\midrule
(a1) premier tirage seul (14) & \textbf{\makecell{+47\\{}[+18\,;\,+65]}} & \textbf{\makecell{+35\\{}[+5\,;\,+54]}} & \textbf{\makecell{+30\\{}[+2\,;\,+54]}} & \textbf{\makecell{+38\\{}[+13\,;\,+52]}} & \makecell{+17\\{}[−8\,;\,+44]} & \makecell{+12\\{}[−1\,;\,+38]} \\
(a1) file A3 seule (36) & \textbf{\makecell{+41\\{}[+19\,;\,+57]}} & \textbf{\makecell{+37\\{}[+16\,;\,+53]}} & \textbf{\makecell{+31\\{}[+10\,;\,+49]}} & \textbf{\makecell{+21\\{}[+1\,;\,+37]}} & \textbf{\makecell{+24\\{}[+4\,;\,+43]}} & \textbf{\makecell{+12\\{}[0\,;\,+27]}} \\
(a1) premier tirage contre A3 : $p$ & 0,73 & 1 & 1 & 0,17 & 0,76 & 1 \\
(a2) arXiv contre site : $p$ & 1 & 1 & 0,15 & 1 & 1 & 1 \\
(a3) Mantel-Haenszel : rapport de cotes & 6,3 & 4,4 & 3,5 & 3,8 & 2,4 & 3,6 \\
\quad $p$ & \textbf{5\,×\,10\textsuperscript{−5}} & \textbf{0,001} & \textbf{0,005} & \textbf{0,006} & 0,06 & 0,15 \\
(a4) un lecteur retiré : écarts & +37 à +50 & +30 à +43 & +27 à +37 & +21 à +31 & +18 à +25 & +9 à +14 \\
(a4) plus basse borne de l'IC & \textbf{+16} & \textbf{+10} & \textbf{+7} & \textbf{+4} & −3 & −2 \\
Fragile selon (a) & non & non & non & non & oui (1) & oui (8) \\
\midrule
Sans les 13 articles opus-5 & \textbf{\makecell{+46\\{}[+26\,;\,+61]}} & \textbf{\makecell{+35\\{}[+16\,;\,+51]}} & \textbf{\makecell{+31\\{}[+11\,;\,+48]}} & \textbf{\makecell{+25\\{}[+8\,;\,+41]}} & \textbf{\makecell{+22\\{}[+3\,;\,+40]}} & \textbf{\makecell{+13\\{}[+1\,;\,+25]}} \\
\quad $p$ Holm & \textbf{0,0001} & \textbf{0,008} & \textbf{0,04} & 0,07 & 0,23 & 0,65 \\
« Non » démentis retournés : $p$ Holm & \textbf{0,001} & \textbf{0,008} & \textbf{0,003} & \textbf{0,04} & 0,24 & 0,66 \\
\quad et douteux aussi : $p$ Holm & \textbf{0,002} & 0,05 & \textbf{0,005} & 0,07 & 0,24 & 0,66 \\
\bottomrule
\end{tabularx}
\end{table}
